% Generated by Listas/make_statements.py from pratica-lecture-08.md, item 14. Do not edit.
\textbf{Source.} Muñoz-Garcia \& Toro-Gonzalez (2019), \emph{Strategy and Game Theory: Practice Exercises with Answers}, Ch. 11, Exercise 5 ``Entry Deterrence Through Price Wars'', p.~505 {[}PDF p.~512{]}. Reproduced verbatim (statement only; the book's worked answer follows it). Answer: yes, the book's worked answer follows the exercise.



\begin{quote}

Consider two firms competing in an industry, an incumbent and an entrant. The entrant has already entered the market, and in period 1 the incumbent must decide whether to fight or accommodate; if she fights, the entrant loses \(f_E > 0\) (i.e., profit \(= -f_E\)) and if she accommodates the entrant makes a profit of \(a_E > 0\). The incumbent may be either ``sane'', in which case he loses \(f_I > 0\) from fighting, and gains \(a_I > 0\) from accommodating; or ``crazy'', in which case his only available action is to fight. In period 2, the entrant must choose whether to stay in the market or exit; if she exits, the incumbent makes a profit of m; if she stays, the incumbent gets to fight or accommodate again and the payoffs are as described before. Assume that there is no discounting, and that the prior probability of a ``sane'' incumbent is \(0 < p < 1\) (the actual type of the incumbent is private knowledge to the incumbent).



\textbf{Part (a)} Draw the game in its extensive form.



\textbf{Part (b)} Show that if payoffs satisfy \(m - f_I \le 2a_I\), then there is a separating equilibrium.



\textbf{Part (c)} Show that if \(m - f_I \le 2a_I\) does not hold there is a pooling equilibrium when \(p \le \frac{f_E}{f_E + a_E}\).

\end{quote}

